\documentclass{amsbook}

\usepackage{enumitem}

\newtheorem{theorem}{Theorem}[chapter]
\newtheorem{lemma}[theorem]{Lemma}

\theoremstyle{definition}
\newtheorem{definition}[theorem]{Definition}
\newtheorem{example}[theorem]{Example}
\newtheorem{xca}[theorem]{Exercise}

\theoremstyle{remark}
\newtheorem{remark}[theorem]{Remark}

\numberwithin{section}{chapter}
\numberwithin{equation}{chapter}

\begin{document}

\frontmatter

\title{Minqi's Solutions to the Exercises of Russell and Norvig (2009), Artificial Intelligence: A Modern Approach (3rd Edition)}

\author{Minqi Pan\footnote{Website: http://www.minqi-pan.com}}

\date{\today}

\maketitle

\setcounter{page}{4}

\tableofcontents

\mainmatter

\chapter{Introduction}

\section{Exercise 1.1}
Define in your own words:
\begin{enumerate}[label=(\alph*)]
\item intelligence,
\item artificial intelligence,
\item agent,
\item rationality,
\item logical reasoning.
\end{enumerate}
\begin{proof}[Answer]
\begin{enumerate}[label=(\alph*)]
\item Intelligence is the ability to perceive information, perform reasoning and thinkings, retain knowledge, and take rational actions.
\item Artificial intelligence is the study of the design of intelligent agents.
\item An agent is an independent entity that operates autonomously, perceives its environment, persists knowledge over a prolonged time period, adapts to changes, and creates and purses goals.
\item Rationality is a characteristic of an agent's behavior, that is, to always achieve the best outcome or the best expected outcome.
\item Logical reasoning is to reason according to the formal rules and norms of logic.
\end{enumerate}
\end{proof}

\section{Exercise 1.2}
Read Turing's original paper on AI (Turing, 1950). In the paper, he discusses several objections to his proposed enterprise and his test for intelligence. Which objections still carry weight? Are his refutations valid? Can you think of new objections arising from developments since he wrote the paper? In the paper, he predicts that, by the year 2000, a computer will have a 30\% chance of passing a five-minute Turing Test with an unskilled interrogator. What chance do you think a computer would have today? In another 50 years?
\begin{proof}[Answer]
Let's examine each of the 9 putative objections discussed by Turing.
\begin{enumerate}
\item The Theological Objection. This objection still carries weight because contemporary knowledge about ``soul'', ``mind'' or ``consciousness'' is still not adequately available to easily trump theological arguments. Turing's refutations are valid, as the teachings of different religions do vary arbitrarily and future scientific development might indeed render theological arguments futile.
\item The ``Heads in the Sand'' Objection. This objection does not carry weight because it is a fallacious ``appeal to consequences'' argument. Turing's refutations are valid, as this objections is not sufficiently substantial.
\item The Mathematical Objection. This objection still carries weight because the implications of Gödel's incompleteness theorem into how intelligence works is still unclear (e.g. some evidence has been produced in recent years regarding the orchestrated objective reduction). Turing's refutations are valid, because whether or not human intelligence is subject to the same limits as machines is a question yet to be answered, and the objection fails in the sense of triumphing simultaneously over all machines.
\item The Argument from Consciousness. This objection still carries weight because it concerns how to achieve artificial consciousness, which is not solved yet. Turing, on the other hand, is only concerned with the external behavior of the machine without considering the internal mental achievements. Turing's refutation is invalid, because he did not directly address the argument by unraveling the mystery of consciousness but instead chose to avoid the argument by retreating to behaviorism. The argument from consciousness and Turing's refutations are based on different standards.
\item Arguments from Various Disabilities. This objection does not carry weight because no support is usually offered for the various disability statements. Turing's refutations are valid, as the disability arguments are mostly founded on the principle of scientific induction yet the versatility of machines is subject to changes in the future, or are ``disguised forms of the argument from consciousness''.
\item Lady Lovelace's Objection. This objection does not carry weight because it is only a historic remark made when evidence available to Lady Lovelace did not encourage her to believe that machines are capable of originality. Moreover, Lady Lovelace's words are not a real objection because it does not preclude the possibility of machine learning (i.e. not using the word ``only'' before ``whatever we know how to order it to perform''). Contemporary development in machine learning has proven it feasible for machines to produce surprising results more than ``whatever we know how to order it'' to perform. Turing's refutations are valid because machines can indeed surprise us and there is of course virtue in the working out of consequences from data and general principles.
\item Argument from Continuity in the Nervous System. This objection does not carry weight because a continuous nervous system is just one way of creating intelligence. This argument reverses necessity and sufficiency and is therefore a fallacy. Turing's refutation is valid, because from his behavioristic point of view it does not make any difference whether the implementation of intelligence is based on a discrete machine or continuous machine. Also it is true that a continuous machine can be simulated by a discrete machine to a reasonable degree of accuracy.
\item The Argument from Informality of Behavior. This objection does not carry weight because it is not necessarily true that the behavior is predictable once the laws/rules governing it have been understood. Moreover, such laws/rules are hardly captured by exhaustion via scientific observations. Turing's argument is valid because he attacked the argument by showing that it was confusing ``laws of behavior'' with ``rules of conduct'' and showed that ``laws of behavior'' could be hardly exhausted and did not lead to feasible behavior predictions.
\item The Argument from Extra-Sensory Perception. This objection does not carry weight because the scientific consensus does not view extrasensory perception as a scientific phenomenon. Turing did not offer refutation in this paper, and instead considers this argument ``quite a strong one''.
\end{enumerate}

Below are devised new objections arising from new developments.
\begin{enumerate}
\item The quantum mind argument. Classical mechanics cannot explain consciousness. Quantum mechanical phenomena may play an important part in the brain's function and could form the basis for an explanation of consciousness. Thus classical machines cannot think.
\item The naturalistic dualism argument. Naturalistic dualists argue that, although mental states are caused by physical systems, they are ontologically distinct from and not reducible to physical systems. Therefore one cannot expect to create mental states from a physical machine.
\item The Chinese room argument. A digital computer running programs is just a computation model of consciousness, instead of the consciousness itself. Therefore such a machine can never achieve true ``understanding'' or true ``consciousness''.
\item The unconscious argument. Human intelligence and expertise depend primarily on unconscious processes rather than conscious symbolic manipulation, and that these unconscious skills can never be fully captured in formal rules. Thus machines that reduce intelligence to symbol manipulation will never have intelligence.
\item The epistemological argument. Much of human knowledge is not symbolic. Therefore a symbol processing machine can never formalize or represent all knowledge to enable intelligence.
\item The ontological argument. It is doubtful whether the world consists of independent facts that can be represented by independent symbols. If this was false, it would be doubtful about what we can ultimately know. Let alone machines.
\item The contextualism argument. Human intelligence is highly context-bound. A machine cannot have human-like intelligence without having a human-like being-in-the-world or having bodies more or less like human beings', or social acculturation (i.e. a society) more or less like human beings'.
\end{enumerate}

I think a computer today has a 10\% chance of passing a five-minute Turing Test with an unskilled interrogator. In another 50 years a computer tends to have 50\% chance of passing a five-minute Turing Test as machines at that time might behave sufficiently deceptive to leave an unskilled interrogator no better than random guessing.
\end{proof}

\section{Exercise 1.3}
Are reflex actions (such as flinching from a hot stove) rational? Are they intelligent?
\begin{proof}[Answer]
They are rational. Because it achieves the best expected outcome according to the result of a lengthy period of evolution. For example, A human flinching from a hot stove has the best chance of protecting the person from being burnt. That immediate flinching is the best course of action is asserted by natural selection and evolution result inherited from the countless ancestors of this person.

They are intelligent. Because it is clear that, a reflex action perceives the external environmental information, performs reasoning (although being a simple and fast reasoning without going through the brain rather than a slow and careful deliberation that goes through the brain), and consequently takes a rational action that produces the best expected outcome.
\end{proof}
\end{document}
